\section{Endpoint means from polynomial vacuum estimates}\label{c:pvmeans}
We now estimate mean lengths for bridges to prescribed top ports.
The polynomial-vacuum representation controls endpoint variation after
summing over heights. This gives individual-height estimates on one
density-one set, and estimates at every large scale after critical
averaging over even heights.

Let $\mathcal B_h(e)$ be the strict bridges from bottom coordinate $0$
to top coordinate $e\in\mathbb Z+h/2$ at height $h$. Define
\[
 b_h(e)=\sum_{\gamma\in\mathcal B_h(e)}\kappa^{L(\gamma)},\qquad
 b_h^{(L)}(e)=\sum_{\gamma\in\mathcal B_h(e)}L(\gamma)\kappa^{L(\gamma)}.
\]
Thus $b_h=\sum_e b_h(e)=B_h$ and $b_0(e)=\mathbf1_{e=0}$.
\begin{theorem}[One set of density-one heights]\label{c:density-one}
There is one set $G\subseteq\mathbb N$ with
$|G\cap[1,N]|/N\to1$ such that, for every fixed $\delta,\varepsilon>0$,
all sufficiently large $h\in G$ satisfy
\[
 h^{4/3-\varepsilon}\le\frac{b_h^{(L)}(e)}{b_h(e)}
       \le h^{4/3+\varepsilon}
 \quad(e\in\mathbb Z+h/2,\ |e|\le h^{1-\delta}).
\]
\end{theorem}
For $E_H=2\mathbb N\cap[H,2H]$, put
$Z_H=\sum_{h\in E_H}b_h(0)$ and
$M_H=\sum_{h\in E_H}b_h^{(L)}(0)$.
\begin{theorem}[All-even aligned height mixture]\label{c:mixture}
Give each pair $(h,\gamma)$, $h\in E_H$ and
$\gamma\in\mathcal B_h(0)$, probability $\kappa^{L(\gamma)}/Z_H$.
As $H\to\infty$ through all real scales,
\[
 Z_H=H^{-1/4+o(1)},\quad M_H=H^{13/12+o(1)},\quad
 \mathbb E_H L=M_H/Z_H=H^{4/3+o(1)}.
\]
The height marginal is $b_h(0)/Z_H$.
\end{theorem}

The complete finite construction in \cite[Theorem 1.1 and Lemma 8.2]{PolynomialVacuum} gives
$b_h=h^{-1/4+o(1)}$ and
$\sum_e b_h^{(L)}(e)=h^{13/12+o(1)}$, with exponential strip aspect
cutoff. The flat-strip flux identity in \cite[Section 2]{PolynomialVacuum}
is $b_h+\cos(3\pi/8)A_h=1$, where $A_h$ is the raw critical
mass of nontrivial paths that stay inside the height-$h$ strip,
from the fixed bottom source back to the bottom wall,
with the terminal port summed. Let $i$ be its positive-height irreducible
probability, now recording height and transverse displacement. Its array
$i^{(L)}(h,e)$ is the sum of $L(\gamma)\kappa^{L(\gamma)}$ over
irreducibles with those height and displacement coordinates. Convolution
below is in both variables. We use the following exact reward input from
\cite[Proposition 10.1]{Renewal}.

\begin{proposition}[Polynomial-vacuum rewards]\label{c:pv-rewards}
Write $J$ for the height of one piece under $i$, and $L$ for its length.
Set $\alpha=1/4$, $\beta=3/4$, and $d_*=4/3$. Then
\[
 b=\sum_{k\ge0}i^{*k},\qquad b^{(L)}=b*i^{(L)}*b,
\]
\[
 i(J>r)=r^{-\beta+o(1)},\quad
 \mathbb E_i[L;J\le r]\le r^{d_*-\beta+o(1)},\quad
 \mathbb E_i[Le^{-J/r}]=r^{d_*-\beta+o(1)}.
\]
The discounted lower reward is retained on pieces with height at most
$r\log^2r$ and lateral excursion at most $r\log^4r$.
\end{proposition}
Only in this section, $J$ denotes height. We use $\sigma=3/8$ in the winding phases and
$\theta=\pi/8$ in their coefficients.

Reflection makes $b_h(e)$ symmetric in $e$. Heights and offsets both add under concatenation, so the two convolution identities apply to these arrays.
The lower estimates below will hold at every large height. The endpoint
upper estimates will instead be summed over heights; this distinction is
what leads to the two different conclusions of the section.

\subsection{Endpoint variation from two disjoint bridges}
First put \(t_h\) equal to the total critical product mass for two vertex-disjoint bridges with adjacent starts 0,1 in that order. Write \(J_h(u,v)\) for its refinement by ordered endpoints \(u<v\) on top. We claim
\[
 b_h(y)=\sum_{u\le y<v}J_h(u,v).
\]
Indeed take a bridge spectator from 1 to \(y'\) on top and grow an active exit from 0 with phase \(e^{i\sigma W}\), stopping instead just before first contact with the spectator vertices. Take also the balance reversing the roles of 0 and 1 (spectator now from 0 to \(y'\)), multiplied by \(s_0=e^{6i\theta}\), \(\theta=\pi/8\). Revisit pairs cancel; convergence at fixed width in the vertical coordinate (the unbounded horizontal strip) follows by the finite strip exponential lateral cutoff.

At a first contact the three arms from \(0,1,y'\) form a tripod with the central weight included once, with bijections from the contact terms of each balance. The three branches have the same turn phases at the first contact, so their turn phases measured before turning at the center satisfy
\(e^{i\sigma W_{y'}}=e^{i\theta}e^{i\sigma W_0}=e^{-i\theta}e^{i\sigma W_1}\).
Hence the contact terms cancel in this combination (\(s_0=-e^{-2i\theta}\)).

The remaining active arches in the first balance can only go left, separated by the bridge from exits beyond the adjacent port (also no exit at it); in the second only right. Their phases including the combination factor are thus both \(\ell=e^{3i\theta}\). Project by \(X\mapsto \Im(\ell^{-1}X)/\Im(\ell^{-1})\), killing them and giving coefficients \(1,-1\) on the two types of bridge terms. This yields
\[
 b_h(y'-1)-b_h(y')=\sum_{u<y'}J_h(u,y')-\sum_{v>y'}J_h(y',v).
\]
Sum up to \(y'=y\).

Consequently \(\sup_e b_h(e)\le t_h\) and the sum of absolute first differences of \(b_h(\cdot)\) on its coset is at most \(2t_h\). Also
\(t_h\le C(A_{h+1}-A_h)\): reflect so the adjacent ports are on the top of the old strip, join them by the bounded minimal arch in the added one layer, and translate horizontally to put the resulting bottom arch's starting port at the fixed source (use the ordered ends). A strip-exchanging lattice isometry with any needed offset can be used before this translation. Recovery has bounded multiplicity since the crossing ports of the old top line are exactly the two joined ports, adjacent and prescribed before translation. Thus by flat strip flux, for fixed \(0<c_1<c_2<\infty\),
\[
 \sum_{c_1 H\le j\le c_2 H} t_j \le C b_{\lfloor c_1H\rfloor}=H^{-\alpha+o(1)}\quad\text{(as an upper bound)} .
\]

\subsection{A pointwise lower bound at every height}
The variation bound is available only after summing heights. Nevertheless
it will yield a lower bound at each height by averaging internal renewal
cuts. Here is a uniform pointwise \textbf{lower} bound at all large heights for any fixed \(\delta>0\):
\[
 b_h(e)\ \ge\ h^{-1-\alpha-o(1)} \qquad (|e|\le h^{1-\delta},\ e\in\mathbb Z+h/2).
\]
First take even \(h\). For a bridge \(\gamma\in\mathcal B_h(e)\), let
\(N_h(\gamma)\) count the integers \(j\in[h/8,h/7]\) for which all
three levels \(j,2j,h/2+j\) are single-crossing cuts. Put
\(k=h/2-j\). Splitting at those cuts gives four bridges of heights
\(j,j,k,k\), so, with convolution here only in the offset,
\[
 \sum_{\gamma\in\mathcal B_h(e)}N_h(\gamma)\kappa^{L(\gamma)}
 =\sum_{h/8\le j\le h/7}(b_j*b_j*b_k*b_k)(e).
\]
For all but \(o(h)\) of these internal levels \(j\), the summed
variation estimate gives \(t_j\le h^{-1-\alpha+\delta_0/8}\), where
\(\delta_0=\min(\delta,1/2)\). The external height \(h\) remains
arbitrary. Write \(A=b_j(\cdot)*b_k(\cdot)\). Its total is
\(b_j b_k\), concentrated on \(|x|\le Ch\log^2 h\) up to negligible
tails by the exponential horizontal strip bounds; its sup norm and total
variation on its coset are bounded by \(t_j b_k,2t_j b_k\)
respectively. Thus
\[
 \sum_x A(x)A(e-x)
 =\sum_x A(x)A(x-e)\ \ge\ \|A\|_2^2-2|e|t_j^2 b_k^2,
 \qquad
 \|A\|_2^2\ge h^{-1-o(1)}(b_jb_k)^2.
\]
For good \(j\) the subtracted term is negligible. Summing gives first moment at least \(h^{-4\alpha-o(1)}\).

For the second moment, diagonal terms cost at most \(h^{-4\alpha+o(1)}\) in sum (use \(\sup b_j(\cdot)\le t_j\) on one factor). Off-diagonal terms \(j<j'\) are cut up by the union of both triples of levels, yielding endpoint convolution masses of seven bridges: four of heights comparable to \(h\), including \(j\), and three of heights \(\Delta,2\Delta,\Delta\), \(\Delta=j'-j\). Bound pointwise using the sup at \(j\) and total masses for all others. Since \(\sum_{\Delta\le h}b_\Delta^2 b_{2\Delta}\le h^{1-3\alpha+o(1)}\), summation of \(t_j\) gives
\[
 \sum_{\gamma\in\mathcal B_h(e)}N_h(\gamma)^2\kappa^{L(\gamma)}
 \le h^{1-7\alpha+o(1)}.
\]
Weighted Cauchy--Schwarz therefore gives
\[
 b_h(e)\ge
 \frac{\bigl(\sum_{\gamma\in\mathcal B_h(e)}
 N_h(\gamma)\kappa^{L(\gamma)}\bigr)^2}
 {\sum_{\gamma\in\mathcal B_h(e)}
 N_h(\gamma)^2\kappa^{L(\gamma)}}
 \ge h^{-1-\alpha-o(1)}.
\]
Odd heights follow by a bounded one-layer extension from even height
with the appropriate half-step endpoint shift, using a lower-bound
window with a smaller positive \(\delta\).

\subsection{Inserting length between two bridge kernels}
Write \(K=b*b\) in both height and offset. It follows that
\[
 K_h(e)\ge h^{-2\alpha-o(1)} \qquad (|e|\le h^{1-\delta})
\]
again for admissible offsets. Indeed sum over heights \(h/3\le l\le 2h/3\) and offsets \(|x|\le h^{1-\eta}\) of the first factor, for arbitrarily small \(0<\eta<\delta/2\); both factors use the preceding lower bound with sublinear power windows. Send \(\eta\) down. On the upper side we have
\[
 \sum_{j\le 2H}\sup_x K_j(x)\le H^{2\beta-1+o(1)} .
\]
Order the two factor heights by size, larger \(l\) and smaller \(r\). Apart from zero total height use \(\sup b_l(\cdot)\le t_l\) and \(\sum_{r\le l} b_r\le l^{\beta+o(1)}\) (up to constants); summing dyadically gives the bound since \(\beta-\alpha>0\).

The length insertion is exactly \(b^{(L)}=K*i^{(L)}\). Uniformly in the same windows we obtain the pointwise lower bound
\[
 b_h^{(L)}(e)\ge h^{d_*+\beta-2-o(1)}.
\]
In fact take the middle piece from \(i^{(L)}\) with height \(\le m\log^2 m\) and absolute displacement \(\le m\log^4 m\), \(m=h^{1-\varepsilon}\), for small fixed \(\varepsilon>0\). Its available total mass is at least \(m^{d_*-\beta-o(1)}\), by the exponentially height-discounted first moment and tail truncations in Proposition~\ref{c:pv-rewards}. The remaining \(K\) uses height comparable to \(h\). Since \(m\log^4m\le h^{1-\varepsilon/2}\) for large \(h\), its offset still lies in a sublinear power window for each fixed \(\delta,\varepsilon>0\). Use the lower bound for \(K\) and send \(\varepsilon\) to zero. Also
\[
 \sum_{H\le h\le 2H}\sup_e b_h^{(L)}(e)\le H^{d_*+\beta-1+o(1)}
\]
by the upper \(K\) estimate and the truncated \(i^{(L)}\) upper bound up to height \(2H\).

\subsection{One density-one set and the aligned mixture}
Put $U_h=\sup_e b_h(e)$ and $V_h=\sup_e b_h^{(L)}(e)$, with the
suprema over the actual row coset. For dyadic $H$, choose
$\rho_H\to\infty$ with $\rho_H=H^{o(1)}$ so that the two summed
upper bounds give
\[
 \sum_{H\le h<2H}U_h\le\rho_H H^{-\alpha},\qquad
 \sum_{H\le h<2H}V_h\le\rho_H H^{d_*+\beta-1}.
\]
In each such interval retain precisely the heights satisfying
\[
 U_h\le\rho_H^2 H^{-1-\alpha},\qquad
 V_h\le\rho_H^2 H^{d_*+\beta-2},
\]
and let $G$ be the union of the retained heights. At most
$2H/\rho_H=o(H)$ heights are removed from each interval. Summing
over dyadic intervals shows that $G$ has density one. On $G$,
\[
 \sup_e b_h(e)\le h^{-5/4+o(1)},\qquad
 \sup_e b_h^{(L)}(e)\le h^{1/12+o(1)}.
\]
The construction of $G$ uses global endpoint suprema and does not
depend on $\delta$ or $\varepsilon$. For each fixed $\delta>0$, the
all-height lower bounds give, uniformly on its admissible window,
$b_h(e)\ge h^{-5/4-o(1)}$ and
$b_h^{(L)}(e)\ge h^{1/12-o(1)}$. Combining the lower numerator with
the upper denominator, and then the upper numerator with the lower
denominator, yields
\[
 h^{4/3-o(1)}\le\frac{b_h^{(L)}(e)}{b_h(e)}
 \le h^{4/3+o(1)}\qquad(h\in G,\ |e|\le h^{1-\delta}).
\]
This proves Theorem~\ref{c:density-one} with its stated quantifiers.

For completeness, the mixture in Theorem~\ref{c:mixture} follows directly from the estimates before that deletion of heights. At each large even height, the aligned lower bounds are $b_h(0)\ge h^{-5/4-o(1)}$ and $b_h^{(L)}(0)\ge h^{1/12-o(1)}$. Summing over the order $H$ even heights in $[H,2H]$ gives $Z_H\ge H^{-1/4-o(1)}$ and $M_H\ge H^{13/12-o(1)}$. The dyadic sums of the endpoint suprema give the reverse bounds $Z_H\le H^{-1/4+o(1)}$ and $M_H\le H^{13/12+o(1)}$. Their ratio proves the mixture mean at every large scale. No upper estimate at an individual discarded height has been used.
